\chapter{Einleitung}
\label{ch:einleitung}

Dieses Kapitel führt in das Thema ein. Es beschreibt zunächst, wie sich
der Einsatz künstlicher Intelligenz auf beiden Seiten eines Angriffs
auswirkt und welche Frage daraus für die Gestaltung von Netzen folgt,
leitet daraus Ziel und Forschungsfrage ab und beschreibt abschließend den
Aufbau der Arbeit.

\section{Motivation und Problemstellung}
\label{sec:motivation}

\ac{KI} ist in der Cybersicherheit auf beiden Seiten
angekommen. Angreifer setzen generative Modelle ein, um Phishing zu
skalieren, Schadcode schneller zu entwickeln und Social Engineering
überzeugender zu machen; die KI-gestützte Angriffsaktivität nahm binnen
eines Jahres um 89~Prozent zu, und vom ersten Zugriff bis zum Übergriff
auf weitere Systeme vergehen im Durchschnitt weniger als dreißig
Minuten~\cite{crowdstrike2026gtr}. Verteidiger nutzen \ac{KI} ihrerseits,
um Telemetrie auszuwerten, Anomalien zu erkennen und Reaktionszeiten zu
verkürzen~\cite{xforce2026}. An den Grundlagen der Angriffe ändert das wenig, denn ausgenutzt werden
weiterhin ungepatchte Lücken, gültige Zugangsdaten und
Fehlkonfigurationen. Verändert haben sich Geschwindigkeit und
Anpassungsfähigkeit. Angreifer, die ihre Taktik in Echtzeit an veränderte
Bedingungen anpassen, setzen Verteidigungen unter Druck, die auf festen
Regeln und Signaturen beruhen~\cite{xforce2026}.

Damit stellt sich die Frage, ob nicht auch die Verteidigung selbst lernen
sollte. \ac{RL} ist dafür ein naheliegender Kandidat: Dabei lernt ein
eigenständig handelndes Programm, ein sogenannter \emph{Agent}, sein
Verhalten aus den Folgen seiner Handlungen statt aus vorgegebenen Regeln und
muss deshalb nicht auf jede Angriffsvariante vorab vorbereitet
werden~\cite{nguyen2023drl}. Mit Simulationsumgebungen wie
Microsofts \ac{CBS}~\cite{cyberbattlesim2021} oder
CybORG~\cite{standen2021cyborg} stehen dafür Trainingsplätze bereit, und
Erweiterungen wie MARLon~\cite{marlon} stellen dem lernenden Angreifer
einen lernenden Verteidiger gegenüber.

Der Blick dieser Forschung richtet sich bislang vor allem auf die Agenten:
auf Algorithmen, Beobachtungsräume und Belohnungsfunktionen. Die Umgebung, in der gelernt wird, bleibt dabei meist eine feste
Randbedingung, das gilt insbesondere für die Struktur des
verteidigten Netzes. In der Praxis gilt
jedoch gerade diese Struktur als tragender Teil der Verteidigung:
Segmentierung bis hin zur Mikrosegmentierung ist ein Grundpfeiler
moderner Sicherheitsarchitekturen~\cite{nist800207}. Ob sich dieser Praxisgrundsatz im Verhalten lernender Verteidiger
niederschlägt, ob ein Agent in einem segmentierten Netz also schneller
lernt und wirksamer verteidigt, ist eine offene Frage. Eine entsprechende Abhängigkeit besteht auf der Gegenseite. Auch der
Angreifer wird auf einer bestimmten Topologie trainiert. Zu prüfen ist daher,
ob seine Wirksamkeit davon abhängt, ob die angegriffene Topologie mit seiner
Trainingstopologie übereinstimmt.

Diese Fragen lassen sich nur in einem Aufbau beantworten, in dem allein
die Topologie variiert, während Inventar, Knotenwerte und
Trainingsbudget gleich bleiben, und in dem die Wirksamkeit an Größen
gemessen wird, die tatsächlich Verteidigungsleistung abbilden. Beides
erwies sich als eigenständige Aufgabe: Der übernommene Aufbau aus MARLon
musste dafür an mehreren Stellen umgebaut werden
(\Cref{ch:versuchsaufbau}).

\section{Zielsetzung und Forschungsfrage}
\label{sec:ziel}

Diese Arbeit untersucht, wie sich die Netzwerktopologie auf
Verteidigungsagenten auswirkt, die mit \ac{RL} trainiert werden. Die Fragestellung gliedert sich in drei Teilfragen: wie schnell ein solcher Agent ein stabiles Verhalten
erreicht, wie wirksam er nach abgeschlossenem Training den Angriff
eindämmt, und welchen Einfluss darauf die Topologie hat, auf der der
Angreifer trainiert wurde. Die zentrale Forschungsfrage lautet:
\begin{quote}
  \textit{Wie beeinflusst die Topologie des verteidigten Netzes das
  Training und die Wirksamkeit eines mit Reinforcement Learning
  trainierten Verteidigungsagenten, und welchen Anteil an diesem Einfluss
  hat die Topologie, auf der der Angreifer trainiert wurde?}
\end{quote}
Der zur Beantwortung entworfene Versuchsaufbau ist Gegenstand von
\Cref{ch:versuchsaufbau}.

\section{Aufbau der Arbeit}
\label{sec:aufbau}

\Cref{ch:grundlagen} legt die Grundlagen. Beschrieben werden das \ac{RL}
und das eingesetzte Verfahren \ac{PPO}, die Grundbegriffe der
Netzwerksicherheit samt den Prinzipien der Netzsegmentierung sowie der Stand
der Forschung zu Simulationsumgebungen für Netzwerkangriffe, aus dem die Wahl
von \ac{CBS} mit der MARLon-Erweiterung hervorgeht.

\Cref{ch:versuchsaufbau} beschreibt den Versuchsaufbau. Dazu gehören die vier
verglichenen Topologien, der Umbau von Aktionsraum und Belohnungsfunktion
gegenüber dem übernommenen Rahmenwerk, die beiden Trainingsphasen, die
erhobenen Zielgrößen und das Konvergenzkriterium sowie der Ablauf der
abschließenden Evaluation.

\Cref{ch:ergebnisse} stellt die Ergebnisse dar und ordnet sie unmittelbar
ein. Auf ein eigenes Diskussionskapitel wird bewusst verzichtet, damit Befund,
Deutung und die zugehörigen Tabellen und Abbildungen an derselben Stelle
stehen. Ein abschließender Abschnitt benennt die Grenzen der Aussagekraft.

\Cref{ch:fazit} beantwortet die Forschungsfrage entlang der drei Teilfragen,
leitet daraus eine Empfehlung für die Praxis ab und benennt Ansatzpunkte für
Anschlussarbeiten. Der Anhang enthält die Hyperparameter sowie die
vollständigen Verläufe aller 24 Kombinationen, die \Cref{ch:ergebnisse} nur
ausschnittweise wiedergibt.
